\section{Heat kernel expansion on manifolds with boundaries}\label{sec5}
\subsection{Two particular cases}\label{s5cases}
We start our analysis of manifolds with boundaries with two simple
examples. First, let us consider a one-dimensional manifold
$M=[0,\pi ]$. Let $D=-\partial_x^2$. We consider both Dirichlet
(\ref{Dirbc}) and Neumann (\ref{Neubc}) boundary conditions (taking
$S=0$ in (\ref{Neubc}) to simplify the calculations). The eigenfunctions
of $D$ are:
\begin{eqnarray}
{\mbox{Dirichlet:}}& &\ \sin (lx),\ l=1,2,\dots ;\nonumber\\
{\mbox{Neumann:}}& &\ \cos (lx), \ l=0,1,2,\dots 
\label{1dimeigen}
\end{eqnarray}
The eigenvalues are $l^2$ in both cases, where $l$ is a positive integer
for Dirichlet boundary conditions, and $l$ is a nonnegative integer for
Neumann boundary conditions. The heat kernel asymptotics can be calculated
with the help of the Poisson summation formula (\ref{eq:asPoi}):
\begin{eqnarray}
&&K(t,-\partial_x^2,\mathcal{B}^-)=\sum_{l>0} \exp (-tl^2)=
\frac 12 \left( \sqrt{\frac{\pi}t} -1 \right) 
+\mathcal{O} \left( e^{-1/t} \right)
\,,\label{1dimDir}\\
&&K(t,-\partial_x^2,\mathcal{B}^+)=\sum_{l\ge 0} \exp (-tl^2)=
\frac 12 \left( \sqrt{\frac{\pi}t} +1 \right) 
+\mathcal{O} \left( e^{-1/t} \right)
\,,\label{1dimNeu}
\end{eqnarray}
where for the later use we explicitly mention the boundary operators
$\mathcal{B}^-$ and $\mathcal{B}^+$ which define Dirichlet (\ref{Dirbc})
and Neumann (\ref{Neubc}) boundary conditions respectively.

Let us modify a little bit the example above by allowing for a non-zero
$S$ at one of the components of the boundary:
\begin{equation}
\partial_x\phi \vert_{x=0}=0,\qquad
(\partial_x-S)\phi \vert_{x=\pi}=0 \label{1dimRob}
\end{equation}
Note, that at $x=\pi$ the derivative with respect to an {\it inward}
pointing unit normal is $-\partial_x$. The eigenfunctions  
\begin{equation}
\phi_k=\cos (kx) \label{1derob}
\end{equation}
satisfy the boundary conditions (\ref{1dimRob}) at $x=0$.
The spectrum is defined by the condition at $x=\pi$ which reads\footnote{
For $S>0$ also a negative mode $\cosh (kx)$ appears. Here we take
$S\le 0$.}:
\begin{equation}
-k \sin (k\pi)=S\cos (k\pi) \,.\label{1dspeceq}
\end{equation}
In this example we restrict ourselves to the linear order in $S$.
We suppose that $S$ is small and that the spectrum (\ref{1dimeigen})
is only slightly
perturbed: 
\begin{equation}
k_l=l+z_l \,.\label{1dpertspec} 
\end{equation}
The equation (\ref{1dspeceq}) gives in this approximation
\begin{eqnarray}
z_l=-\frac{S}{l\pi}&\qquad\mbox{for}\qquad& l>0\,,\nonumber \\
z_0^2=-\frac{S}{\pi}&\qquad\mbox{for}\qquad& l=0\,.\label{1dpertsp}
\end{eqnarray}
We expand also the heat kernel in a power series in $S$ and use
(\ref{1dpertsp}) and (\ref{1dimDir}) to obtain:
\begin{eqnarray}
K(t,-\partial_x^2,\mathcal{B}^+_S)&=&\sum_{l\ge 0} \exp (-tk_l^2)
\nonumber\\
&=&\sum_{l> 0} \exp (-tl^2) \left( 1+\frac{2St}{\pi} \right)
+1+\frac{St}{\pi} +\mathcal{O}(S^2)\nonumber\\
&=&\frac 12 \left( \sqrt{\frac{\pi}t} +1 \right) +
S\sqrt{\frac t\pi}  +\mathcal{O}(S^2) \,,\label{1dlinS}
\end{eqnarray}
where we also dropped exponentially small terms in $t$.

Our next example is the heat kernel expansion for scalar fields in
a ball with Dirichlet and Neumann boundary conditions
\cite{Stewartson:1971,Waechter:1972,Kennedy:1978jp,Moss:1989yf,Bordag:1996gm}.
The metric of the unit ball in $\mathbb{R}^n$ reads:
\begin{equation}
ds^2=dr^2+r^2 d\Omega^2, \quad 0\le r \le 1 , \label{ballmetric}
\end{equation}
where $d\Omega^2$ is the metric on unit sphere $S^{n-1}$.
The scalar Laplace operator has the form
\begin{equation}
\Delta \phi =
\left( \partial^2_r+\frac {n-1}r \partial_r+^{(n-1)}\Delta 
\right) \phi \,,
\label{ballLapl}
\end{equation}
where $^{(n-1)}\Delta$ is the Laplace operator on $S^{n-1}$.
The eigenfunctions
of the operator (\ref{ballLapl}) are well known
\begin{equation}
\phi_{l,\lambda }\propto r^{(2-n)/2 } J_{(n-2)/2 +l}
(r\lambda^{1/2} ) Y_{(l)}(x^a). \label{balleigenf}
\end{equation}
$J_p$ are the Bessel functions. The eigenvalues $-\lambda$
are defined by boundary conditions. $Y_{(l)}(x^i)$ are
$n$-dimensional scalar spherical harmonics. Their degeneracies
are
\begin{equation}
N_l=\frac {(2l+n-2)(l+n-3)!}{l!(n-2)!} . \label{spherdeg}
\end{equation}

To proceed further it is convenient to consider the zeta function
which may be presented through a contour integral 
\cite{Bordag:1996gm}\footnote{A general discussion of the representation
of the zeta function by contour integrals can be found in 
\cite{Kirsten:2003py}.}:
\begin{equation}
\zeta (s,D)=\sum_{l=0}^\infty N_l \int_\gamma
\frac{dk}{2\pi i} k^{-2s} \frac{\partial}{\partial k}
\Phi_{(n-2)/2 +l} (k) \,,\label{ballzeta}
\end{equation}
where $\Phi_\nu$ is a function which has zeros at the spectrum
$k=\sqrt\lambda$.
For Dirichlet boundary conditions this function reads:
\begin{equation}
\Phi_\nu =k^{-\nu} J_\nu (k) \,.\label{ballPhi}
\end{equation}
For Robin boundary condition $\Phi_\nu$ is given by a somewhat
more complicated combination of the Bessel functions. The contour 
$\gamma$ runs counterclockwise and encloses all the solutions
of  $\Phi_\nu =0$ on the positive real axis. Note the presence 
of $k^{-\nu}$ in (\ref{ballPhi}) which is included to avoid unwanted
contributions coming from the origin $k=0$.

The next step is to rotate the contour to the imaginary axis and
to calculate residues of $\Gamma (s)\zeta (s,D)$ (as prescribed
by (\ref{aRes})). It is interesting to note that the heat kernel
coefficients are defined by several leading terms in the uniform
asymptotic expansion of $\Phi_\nu (i\nu z)$ for large $\nu$ and
fixed $z$. For further details we refer to  \cite{Bordag:1996gm}.

Euclidean ball was frequently used in calculations of the heat kernel
coefficients and functional determinants. Apart from the papers
already quoted above also the computations for a scalar field 
\cite{Dowker:1995sp,Bordag:1996zc,Dowker:1996bh,
Bordag:1997ma,Nesterenko:1999bb}, spinors 
\cite{D'Eath:1991td,Kamenshchik:1993vj,Kamenshchik:1994tw,Dowker:1996ny,Dowker:1996sw,Kirsten:1996sy,Elizalde:1998hx,Esposito:2002vz},
abelian gauge fields 
\cite{Louko:1988ay,Esposito:1994bv,Esposito:1994xe,Vassilevich:1995we,Vassilevich:1995zn,Dowker:1996ny,Kirsten:1996sy,Esposito:1996hu,Esposito:1996gw,Elizalde:1997zw,Esposito:2000dy,Bernasconi:2003}\footnote{Some of these works consider contributions of the
so called physical modes only. As explained in \cite{Vassilevich:1995cz}
complete answer for the effective action and for the scaling behaviour
must include contributions from ghosts and non-physical modes. This
applies also to the spherical cap case considered below.}
should be mentioned.
A similar technique works also for more complicated
geometries as, e.g., the spherical cap 
\cite{Barvinsky:1992tr,Kamenshchik:1992gz,Barvinsky:1992dz,Kamenshchik:1993vj,
Dowker:1995sp}.
Note that here we consider local boundary conditions without tangential
derivatives only. For other types of boundary operators see sec.\
\ref{s5other}.
%%%%%%%%%%%%%%%
\subsection{Dirichlet and Neumann boundary conditions}\label{s5DN}
Let us now find analytic expressions for the heat kernel coefficients
in terms of geometric invariants. Here we follow the method of
Branson and Gilkey \cite{Branson:1990a,Gilkey:1995}. For both 
Dirichlet and modified Neumann (Robin) boundary conditions the
heat kernel coefficients are locally computable. This means
that $a_k$ may be represented as a sum of volume and boundary
integrals of some local invariants:
\begin{equation}
a_k(f,D,\mathcal{B})=\iM f(x) a_k(x,D)
+\sum_{j=0}^{k-1} \idM f^{(j)} a_{k,j}(x,D,\mathcal{B})\,,
\label{akb-loc}
\end{equation}
where $f^{(j)}$ denotes $j$-th normal derivative of the smearing
function $f$. Note, that we cannot now integrate by parts to remove
normal derivatives from $f$. This reflects the distributional nature
of the heat kernel asymptotics.

The volume terms $a_k(x,D)$ are the same as in the previous section
(see (\ref{a0nobou}) - (\ref{a6nobou})). Here we evaluate the
boundary terms $a_{k,j}(x,D,\mathcal{B})$. Obviously, the canonical
mass dimension of $a_{k,j}$ is $k-j-1$. In addition to the usual bulk
invariants $E$, $R$, etc, $a_{k,j}$ can also contain specific boundary
quantities as the extrinsic curvature $L_{ab}$ or $S$ (for Robin
boundary conditions). Note, that $L_{ab}$ and $S$ are defined on the
boundary only and, therefore, can be differentiated only 
tangentially\footnote{See sec.\ \ref{s2geo} for more information on
differential geometry of manifolds with boundary.}. Canonical mass dimension
of $L_{ab}$ and $S$ is $+1$.

In this section we explicitly calculate the first three heat kernel
coefficients $a_0$, $a_1$ and $a_2$. Basing on the considerations of the
preceding paragraph we may write:
\begin{eqnarray}
&&a_0(f,D,\mathcal{ B}^\pm)=(4\pi)^{-n/2}\iM {\rm tr}_V(f),\label{a0DN} \\
&&a_1(f,D,\mathcal{ B}^\pm)=(4\pi)^{-(n-1)/2}\idM \, b_1^\pm
     {\rm tr}_V ( f), \label{a1DN} \\
&&a_2(f,D,\mathcal{ B}^\pm)=\frac 16 (4\pi)^{-n/2}\left\{ \iM
     {\rm tr}_V(6fE+fR) \right. \nonumber \\
&&\qquad\qquad \left. +\idM {\rm tr}_V(b_2^\pm fL_{ aa}
     +b_3^\pm f_{ ;n}+b_4^\pm fS)
     \right\} ,\label{a2DN}
\end{eqnarray}   
where $\mathcal{B}^\pm$ denotes either Robin ($\mathcal{B}^+$) or
Dirichlet ($\mathcal{B}^-$) boundary operator. $b^\pm$ are some
constants. To keep uniform notations we formally included $S$ also in the
expression for Dirichlet boundary conditions. It should be assumed
that $S=0$ for that case.

Since the product 
formula (\ref{product}) is still valid with obvious modifications
for the boundary contributions, we can again consider the
case $M=M_1\times S^1$ and repeat step by step the calculations
(\ref{77hk}) - (\ref{uIn}) thus arriving to the same conclusion that
the constants $b^\pm$ do not depend on dimension of the manifold
(all explicit $n$-dependence of the heat kernel is contained in the power
of $4\pi$). This property is crucial for our calculations.

The one dimensional examples\footnote{For $M=[0,\pi ]$ the boundary
integral reduces to a sum of two contributions from $x=0$ and $x=\pi$.}
(\ref{1dimDir}) and (\ref{1dimNeu}) immediately give
\begin{equation}
b_1^+=-b_1^-=\frac 14 \,.\label{b1pm}
\end{equation}
Our next example (\ref{1dlinS}) controls linear terms in $S$ and
gives
\begin{equation}
b_4^+=12 \,.\label{b4p}
\end{equation}
For Dirichlet boundary conditions we have set $S=0$. Therefore,
$b_4^-$ plays no role.

To define the constants $b_2^\pm$ and $b_3^\pm$ we use the conformal
variation equation (\ref{vara}). Variations of $E$ and $R$ are given
by (\ref{cvar1}). Conformal properties of the second fundamental
form $L_{ab}$ are dictated by that of the metric:
\begin{equation}
\frac d{d\epsilon} \vert_{\epsilon =0} L_{aa}=-fL_{aa} 
-(n-1)f_{;n} \,.\label{cvarLaa}
\end{equation}
We wish to keep the boundary conditions invariant under the conformal
transformations. This implies that the boundary operators $\mathcal{B}^\pm$
must transform homogeneously. This property holds automatically for
Dirichlet boundary conditions (\ref{Dirbc}). In the case of modified
Neumann boundary conditions conformal transformation of $S$ should
cancel inhomogeneous term in the connection $\omega_n$. This yields
\begin{equation}
\frac d{d\epsilon} \vert_{\epsilon =0}S=-fS+\frac 12 (n-2)f_{;n}\,.
\label{cvarS}
\end{equation}
Next we substitute the variational formulae (\ref{cvar1}), (\ref{cvarLaa})
and (\ref{cvarS}) in (\ref{vara}) with $k=2$. We collect the
terms with $f_{;n}$ on the boundary to obtain
\begin{equation}
-(n-4)-(n-1)b_2^-=(n-2) b_3^- \label{cvara2b}
\end{equation}
for Dirichlet conditions and
\begin{equation}
-(n-4)-(n-1)b_2^+ +6(n-2)=(n-2) b_3^+ \label{cvara2N}
\end{equation}
for generalised Neumann (Robin) ones. Since the constants $b^\pm$ do
not depend on $n$ the two equations (\ref{cvara2b}) and (\ref{cvara2N})
are enough to define these constants:
\begin{equation}
b_2^+=b_2^-=2\,,\qquad b_3^+=-b_3^-=3 \,.\label{bpm23}
\end{equation}
This completes the calculation of $a_2$. 

The coefficients $a_3$ and $a_4$ can be obtained as particular
cases of more general formulae of the next subsection.
%%%%%%%%%%%%%%
\subsection{Mixed boundary conditions}\label{s5mixed}
Let us now turn to mixed boundary conditions (\ref{mixbc}) which,
as we have seen in sections \ref{s3spin} and \ref{s3vect}, are
natural boundary conditions for spinor and vector fields.  
These boundary conditions depend on two complementary local
projectors $\Pi_-$ and $\Pi_+=1-\Pi_-$ which define subsets
of components of the field satisfying Dirichlet and Robin
boundary conditions respectively\footnote{Examples of such projectors
are given by eq.\ (\ref{proBG}) and below eq.\ (\ref{relbc}).}.
More precisely,
\begin{equation}
\Pi_-\phi \oB =0\,,\qquad (\nabla_n+S)\Pi_+\phi \oB =0\,.\label{mixagain}
\end{equation}
Consequently, there is one more independent entity (as compared to Robin
case, for example) on which the heat kernel coefficients for mixed 
boundary conditions can depend. This makes the calculations somewhat
more complicated. It is convenient to define
\begin{equation}
\chi = \Pi_+-\Pi_- \,.\label{defchi}
\end{equation}

Calculation of the coefficients $a_k$ up to $k=4$ can be found
in \cite{Branson:1990a} (see also \cite{Vassilevich:1995we} for
some corrections). The result reads:
\begin{eqnarray}
&&a_0(f,D,\mathcal{ B})=(4\pi)^{-n/2}\iM {\rm tr}_V(f).\label{a0bou} \\
&&a_1(f,D,\mathcal{ B})={\frac 14}(4\pi)^{-(n-1)/2}\idM
     {\rm tr}_V (\chi f). \label{a1bou} \\
&&a_2(f,D,\mathcal{ B})=\frac 16 (4\pi)^{-n/2}\left\{ \iM
     {\rm tr}_V(6fE+fR) \right. \nonumber \\
&&\qquad\qquad \left. +\idM {\rm tr}_V(2fL_{ aa}+3 \chi f_{ ;n}+12fS)
     \right\} .\label{a2bou} \\
&&a_3(f,D,\mathcal{ B})=\frac 1{384}(4 \pi )^{
       -(n-1)/2}  \idM \ptr\big\{ f(96 \chi E+16\chi  R \nonumber\\
&&\qquad\qquad+8f \chi R_{anan} 
       +(13 \Pi_{+}-7 \Pi_{ -})L_{ aa}L_{ bb}+(2 \Pi_{ +}+10
      \Pi_{-})L_{ab}L_{ ab}\nonumber\\
&&\qquad\qquad+96SL_{ aa}
      +192S^2 
       -12 \chi_{ :a} \chi_{:a})+f_{ ;n}
      ((6 \Pi_{ +}+30 \Pi_{ -})L_{
        aa}\nonumber\\
&&\qquad\qquad+96S)+24 \chi f_{ ;nn} \}.\label{a3bou}\\
&&a_4(f,D,\mathcal{ B})=\frac 1{360} (4 \pi )^{
       -n/2} \big\{ \iM \ptr
       \{ f(60E_{ ; ii}+60 R E+180E^2 \nonumber \\
&&\qquad\qquad +30 \Omega_{ij} 
      \Omega_{ij} +12 R_{;ii}+
      5 R^2-2R_{ij}R_{ij}+2R_{ijkl}R_{ijkl}) \} \nonumber \\
&&\qquad\qquad+ \idM \ptr
      \big\{ f \{ (240 \Pi_{ +}-120 \Pi_{ -})E_{ ;n}\nonumber\\
&&\qquad\qquad+(42 \Pi_{
     +}-18 \Pi_{ -}) R_{ ;n} 
     +24L_{ aa:bb}+0L_{ ab:ab}+120EL_{
       aa}\nonumber\\ 
&&\qquad\qquad +20 R L_{ aa}+4R_{ an    an}L_{bb} 
      -12R_{ an    bn}L_{ ab}+4R_{
       ab    cb}L_{ ac}\nonumber\\
&&\qquad\qquad +\frac 1{21} \{ (280 \Pi_{ +}+40 \Pi_{
-})L_{ aa}L_{ bb}L_{ cc} 
    +(168 \Pi_{ +}\nonumber\\
&&\qquad\qquad -264 \Pi_{
-})L_{ ab}L_{ ab}L_{ cc}+(224 \Pi_{
+}+320 \Pi_{ -})L_{ ab}L_{ bc}L_{ac} \} \nonumber \\
&&\qquad\qquad +720SE+120S R +0SR_{ an    an}+144SL_{
aa}L_{ bb}+48SL_{ ab}L_{ ab} \nonumber \\
 &&\qquad\qquad+480S^2L_{ aa}+480S^{
3}+120S_{ :aa}+60 \chi  \chi_{ :a} \Omega_{ an}-12 \chi_{ :a} \chi_{
:a}L_{ bb} \nonumber \\
&&\qquad\qquad -24 \chi_{ :a} \chi_{
:b}L_{ ab}-120 \chi_{ :a} \chi_{
:a}S \} +f_{ ;n}(180 \chi E+30 \chi  R +0R_{
an an} \nonumber \\
&&\qquad\qquad +\frac 17 \{ (84 \Pi_{
+}-180 \Pi_{ -})L_{ aa}L_{ bb}+(84 \Pi_{
+}+60 \Pi_{ -})L_{ ab}L_{ ab} \} \nonumber \\
&&\qquad\qquad +72SL_{ aa}+240S^2-18 \chi_{
:a} \chi_{ :a})+f_{ ;nn}(24L_{
aa}+120S) \nonumber \\
&&\qquad\qquad
\left. +30 \chi f_{ ;iin}  \big\} \right\} . \label{a4bou}
\end{eqnarray}
Dirichlet and modified Neumann (Robin) boundary conditions 
are recovered when $\Pi_+=0$
or $\Pi_-=0$ respectively.

The coefficients $a_k$ with $k=0,1,2,3,4$ were calculated
in  \cite{McKean:1967,Kennedy:1980ar,Smith:1981,Melmed:1988hm,Moss:1989mz,
Branson:1990a,McAvity:1991we,McAvity:1992rf,Vassilevich:1995we}.
A different algorithm was suggested in \cite{Cognola:1990kq}.
For ``pure'' (not mixed) boundary conditions the 
coefficient $a_5$ was calculated
by Branson, Gilkey and Vassilevich \cite{Branson:1997cm} for the
special case of a domain in flat space or of a curved domain with
totally geodesic boundaries. Kirsten \cite{Kirsten:1997qd} generalised
these results for arbitrary manifolds and boundaries. Branson, Gilkey,
Kirsten and Vassilevich \cite{Branson:1999jz} calculated the coefficient
$a_5$ for mixed boundary conditions.
%%%%%%%%%%%%%%
\subsection{Other boundary conditions}\label{s5other}
From the technical point of view boundary conditions for a Laplace
type operator are needed to exclude infinite number of negative and
zero modes and to ensure self-adjontness. In principle, any linear
relation between the boundary data $\phi\oB$ and $\phi_{;n}\oB$
is admissible as long as it serves this purpose. Here we consider
two physically motivated examples of boundary conditions which
contain tangential derivatives on the boundary. These examples should
give the reader an idea of what can be expected for a more general 
boundary value problem.
%%%%%%%
\subsubsection{Boundary conditions with tangential derivatives and
Born--Infeld action from open strings}\label{s5tang}
The boundary condition
\begin{equation}
\left. \left( \nabla_n +\frac 12 ( \nabla_a \Gamma_a + \Gamma_a \nabla_a )
+S \right) \phi\right\oB=0 \label{obliquebc}
\end{equation}
is the simplest condition containing both normal and tangential
derivatives. $\Gamma_a$ and $S$ are some matrix valued functions
defined on the boundary. Such or similar structures appear in
open strings (cf. sec.\ \ref{s3string}) and in chiral bags
(cf. sec.\ \ref{s3spin}). They also describe photons with the
Chern-Simons interaction term concentrated on the boundary
\cite{Elizalde:1998ha,Bordag:1999ux}
and may be relevant for solid state physics applications.
The conditions (\ref{obliquebc}) appeared in the mathematical
literature \cite{Grubb:1974,Gilkey:1983}.

Several heat kernel coefficients for the boundary conditions
(\ref{obliquebc}) (which are called {\it oblique}) have been
calculated by McAvity and Osborn \cite{McAvity:1991xf} and
by Dowker and Kirsten \cite{Dowker:1997mn,Dowker:1998hm}.
Avramidi and Esposito 
\cite{Avramidi:1998sb,Avramidi:1998sh,Avramidi:1997hy,Avramidi:1998xj}
lifted some commutativity assumptions and proved a simple
criterion of strong ellipticity (see below).

Consider a differential operator $Q$, not necessarily of the second order.
Let is separate the part containing highest derivatives and replace
$\partial_\mu \to ik_\mu$, like in doing the Fourier transformation.
In this way we obtain an object $A_Q=a^{\mu\nu\dots\rho}(x)k_\mu k_\nu
\dots k_\rho$ which is called the leading symbol of $Q$. If the operator
$Q$ acts on fields with indices (i.e. in a vector bundle), the index structure
of $Q$ is inherited by $A_Q$. In other words, for fixed $k$ the leading symbol
$A$ is a matrix valued function (an endomorphism of the vector bundle).
If $A_Q(x)$ is non-degenerate for all $k\ne 0$, the operator $Q$ is called
elliptic. For the Dirac operator $Q=\Dir$ the leading symbol is 
$A_{Q} =\gamma^\mu k_\mu$. Therefore, $\Dir$ is obviously elliptic.
The Laplace type operators are also elliptic since $A_D=k^2$. 
Ellipticity means that ``at large momenta'' the operator $Q$ is dominated
by its' highest derivative part. This is the property which guarantees
that on a compact manifold  without boundaries Laplace operators have
at most finite number of negative and zero eigenvalues and which
ensures existence of the heat kernel. On manifolds with boundaries
just ellipticity is not enough to ensure nice properties of the spectrum.
There is an additional requirement, called {\it strong ellipticity},
which should be satisfied by the boundary operator (see 
\cite{Gilkey:1995} for details). Dirichlet and Robin boundary conditions
are always strongly elliptic. Oblique boundary conditions are strongly
elliptic if and only if $|\Gamma^2|<1$. If this inequality is violated,
infinite number of negative eigenmodes appears (a simple example can be
found in Appendix B of \cite{Kummer:2000ae}). Note that the boundary 
conditions (\ref{obliquebc}) correspond to a symmetric $D$ if 
the matrices $\Gamma_a$ are anti-hermitian. Therefore, 
$\Gamma^2=\Gamma_a\Gamma^a$ is typically negative.

Another complication stems from the fact that $\Gamma_a$ is dimensionless.
Consequently, arbitrary powers of $\Gamma_a$ may enter $a_k$, so that instead
of the undetermined constants (cf. (\ref{a0DN}) - (\ref{a2DN}))
one has to deal with undetermined functions of $\Gamma$. The problem
becomes more tractable if we suppose that $\Gamma_a$'s commute among
themselves, $[\Gamma_a,\Gamma_b]=0$. For this case, the coefficients
$a_1$ and $a_2$ have been calculated by McAvity and Osborn 
\cite{McAvity:1991xf} ($a_0$ is still given by (\ref{a0DN})):
\begin{eqnarray}
&&a_1(f,D)=(4\pi )^{-(n-1)/2} \idM \ptr (f\gamma (\Gamma ))\,,
\label{a1ob}\\
&&a_2(f,D)=\frac 16 (4\pi)^{-n/2}\left\{ \iM
     {\rm tr}_V(6fE+fR) \right. \nonumber \\
&&\qquad\qquad +\idM {\rm tr}_V(b_0(\Gamma) fL_{ aa}
     +b_1(\Gamma) f_{ ;n}\nonumber\\
&&\qquad\qquad \left. +b_2(\Gamma) fS+f\sigma (\Gamma) L_{ab}
     \Gamma_a\Gamma_b )
     \right\} ,\label{a2ob}
\end{eqnarray}   
where
\begin{eqnarray}
&&\gamma =\frac 14 \left[ \frac 2{\sqrt{1+\Gamma^2}} -1\right]\,,
\nonumber\\
&&b_0=6\left[ \frac 1{1+\Gamma^2}-\frac 1{\sqrt{-\Gamma^2}}
{\rm artanh} \left( \sqrt{-\Gamma^2} \right) \right]+2\,,\nonumber\\
&&b_1=\frac 6{\sqrt{-\Gamma^2}} {\rm artanh} \left( \sqrt{-\Gamma^2} \right)
-3 \,,\label{obliqueconst}\\
&&b_2=\frac {12}{1+\Gamma^2} \,,\nonumber\\
&&\sigma = \frac 1{\Gamma^2} (2-b_0) \,.\nonumber
\end{eqnarray}
We see, that if $\Gamma^2$ approaches $-1$ the heat kernel blows up
indicating violation of the strong ellipticity condition.

As an example, let us consider the open string sigma model of sec.\
\ref{s3string}. To simplify the subsequent calculations we put $B_{AB}=0$.
In the zeta function regularization divergent part of the effective
action (\ref{Wzpole}) reads $(1/s)a_2(D)=(1/s)a_2(1,D)$. On general grounds
we expect that $a_2$ repeats the structure of the classical action
(\ref{actsigma}):
\begin{equation}
a_2=\int_{\partial M}d\tau  \beta_A^{[A]} \partial_\tau \bar X^A +
\mbox{bulk terms},\label{divstring} 
\end{equation}
where $\beta_A^{[A]}$ is a beta 
function\footnote{In this case the subscript $A$
is a target space vector index, while the superscript $[A]$ indicates the
coupling. Here we do not consider other beta functions $\beta^{[G]}$,
$\beta^{[B]}$ etc.}. We put $L_{ab}=0$ (otherwise we have had to introduce
a dilaton coupling on the boundary to achieve renormalizability).
We also suppose that the target space metric $G_{AB}$ is trivial.
With these simplifying assumptions the beta function can be easily
calculated from (\ref{divstring}), (\ref{a2ob}), (\ref{obliqueconst}),
(\ref{bopsigma}):
\begin{equation}
\beta^{[A]}_C=-\frac 1{2\pi} (\partial_A F_{BC})
(1+F^2)^{-1}_{BA} \,.\label{betaAC}
\end{equation}
By a lengthy but straightforward calculation one can demonstrate
that the condition $\beta^{[A]}_C=0$ is equivalent to the
equations of motion following from the Born-Infeld action on the
target space:
\begin{equation}
\mathcal{L}_{\rm BI}=\int d\bar X \sqrt{\det (1+iF)}=
\int d\bar X \exp \left( \frac 14 {\rm tr} \ln
(1+F^2) \right)\,,\label{BIaction}
\end{equation}
where $i$ appeared due to our rule of the Euclidean rotation
for the gauge fields. 

The Born-Infeld action has been derived from the beta functions
of the open string sigma model in \cite{Abouelsaood:1987gd,Callan:1987bc}
confirming an earlier work \cite{Fradkin:1985qd} which used different 
methods. The heat kernel analysis was performed by Osborn 
\cite{Osborn:1991gm} and then repeated in \cite{Kummer:2000ae}
for more general couplings. 
%%%%%%
\subsubsection{Spectral or Atiyah--Patodi--Singer (APS) boundary conditions}
\label{s5spec}
Spectral boundary condition were introduced by Atiyah, Patodi and
Singer in their study of the Index Theorem 
\cite{Atiyah:1975jf,Atiyah:1976jg,Atiyah:1980jh}. These boundary conditions
are {\it global}, i.e. they cannot be defined by using
local data only.

Consider a Dirac type operator 
\begin{equation}
\Dir =i \gamma^\mu \tilde \nabla_\mu +\mathcal{E}\,,\label{gDir}
\end{equation}
where $\tilde\nabla$ is a covariant derivative with a compatible
connection:
\begin{equation}
\tilde \nabla_\mu \gamma_\nu =0,\label{compaticon}
\end{equation}
i.e., the gamma matrices are covariantly constant. $\mathcal{E}$
is a zeroth order operator (a matrix valued function). We also 
suppose that the connection in $\tilde\nabla$ is unitary. This means
that the connection one-form is represented by an anti-hermitian
matrix in a suitable basis. We restrict ourselves to the case
when $\mathcal{E}^\dag =\mathcal{E}$, so that the operator
$\Dir$ is formally self-adjoint in the bulk. Note, that compatible
unitary connection is not unique. 
Consider a first
order differential operator on the boundary:
\begin{equation}
P=\gamma_n \gamma^a \tilde\nabla_a +\frac i2 \left(
\mathcal{E}\gamma_n -\gamma_n \mathcal{E} \right) +\Theta (x)\,,
\label{Poper}
\end{equation}
where $\Theta (x)$ is a hermitian matrix valued function on $\partial M$.
The operator $P$ is a self-adjoint operator of Dirac type on the boundary. 
All functions on $\partial M$ can be decomposed in positive, negative,
and zero modes of $P$. Let us define $\Pi_-$ as a projector on the
space spanned by {\it non-negative} eigenspaces of $P$. Then the
equation
\begin{equation}
\Pi_-\phi \oB=0 \label{APSPmin}
\end{equation}
defines the APS boundary conditions. 

Of course, relative complexity and non-locality limits physical
applications of spectral boundary conditions. However, they appeared
in a number of axial anomaly calculations 
\cite{Romer:1977qq,Hortacsu:1980kv,Hortacsu:1983iq,
Ninomiya:1985ge,Niemi:1986ht,Ma:1986qp,
Forgacs:1987ti,Falomir:1998as}, quantum cosmology \cite{D'Eath:1991sz},
and in works on the Aharonov-Bohm effect
\cite{Mishchenko:1992vm,Beneventano:1998ai}.
In brane models spectral boundary conditions describe T-selfdual
configurations which may be interpreted as mixtures of D-branes
and open strings \cite{Vassilevich:2001at}.

General form of the heat kernel expansion for spectral boundary
conditions was established by
Grubb and Seeley \cite{GS:1993,GS:1995,GS:1996}:
\begin{equation}
K(t,f,D)\simeq \sum_{k=0}^{n-1} a_kt^{(k-n)/2}
+\sum_{j=n}^\infty \left( a'_j \ln t +a''_j \right)
t^{(j-n)/2} \,.\label{APSasymp}
\end{equation}

In the contrast to all previous cases the expansion (\ref{APSasymp}) contains
logarithms of the proper time $t$. Although, such terms appear ``typically''
\cite{GG98} for pseudo-differential operators, they may lead to rather
unpleasant physical consequences. As follows from (\ref{heatzeta}),
non-zero $a'_n$ means that the zeta function has a pole at $s=0$, and,
therefore, the expression (\ref{Wzren}) for the renormalised
effective action does not make sense. Fortunately, for the APS
boundary conditions $a'_n=0$ if $f=1$ near the boundary. 
Consequently, the integrated
zeta function is regular at $s=0$ and eq.\ (\ref{Wzren}) still can be
used. However, all calculations involving localised heat kernel
coefficients remain problematic. Many logarithmic terms vanish if
the manifold $M$ has a product structure near $\partial M$.
In non-product cases strong criteria of ``partial vanishing of
logarithms'' have been found recently \cite{Grubb:2003ys}.

Another problem with the heat kernel expansion (\ref{APSasymp})
is that the coefficients may have a more complicated dependence
on $n$ than just a power\footnote{Let us remind that
to prove the simple dependence on $n$ we used the product
formulae (\ref{77hk}), (\ref{eq:asPoi}). We assumed that the
spectral problem can be ``trivialised'' in one direction.
In the present case all tangential coordinates enter the operator
(\ref{Poper}) on equal footings. Therefore, the proof does not go through.}
 of $(4\pi )$. After this long list of troubles it is not a surprise
that {\it only} the coefficients $a_k$ with $k<n$ are locally
computable.
The coefficient $a_0$ is given by (\ref{a0nobou}).
$a_1$ and  $a_2$ have been calculated in \cite{Dowker:2000sy}:
\begin{eqnarray}
&&a_1=(4\pi )^{(1-n)/2} \frac 14 (\beta (n)-1) \idM \ptr (f)\,,
\nonumber\\
&&a_2=(4\pi )^{-n/2}\left[ \iM \ptr \left( f\left(\frac 16 R+E\right) \right)
\right.\nonumber\\
&&\qquad\qquad +\idM \ptr \left( f \mathcal{E} +\frac 13 \left( 1-\frac 34 
\pi \beta (n) \right)L_{aa}f \right. \label{hkAPS} \\
&&\qquad\qquad \left.\left. -\frac{n-1}{2(n-2)} \left( 1-\frac 12 \pi
\beta (n)\right) f_{;n} \right) \right] \,,\nonumber
\end{eqnarray}
where
\begin{equation}
\beta (n)=\Gamma \left (\frac n2 \right) \Gamma \left( \frac 12 \right)^{-1}
\Gamma \left( \frac{n+1}2 \right)^{-1} \,.\label{betaAPS}
\end{equation}
The coefficient  $a_3$ can be found in \cite{Gilkey:2000bc}.

Some string theory applications suggest 
\cite{Vassilevich:2001at,Vassilevich:2001wf} that spectral boundary conditions
can be defined directly for a second order differential operator.
Existence of the asymptotic expansion (\ref{APSasymp}) and vanishing of
leading logarithms for such problems have been stated in
\cite{Grubb:2003yr}.